\begin{center}
    {\LARGE \textbf{2019物理化学}} \\[2ex]
    {\large 时量180分钟 \quad 满分150分}
\end{center}

\vspace{0.5cm}
\noindent\textbf{注意事项:}
\begin{enumerate}[label=\arabic*., parsep=0pt, itemsep=0pt]
    \item 答题前,考生先将自己的姓名、学号、考场号、座位号填写在答题卡上,将条形码准确粘贴在考生信息条形码粘贴区。
    \item 考试结束后,将本试卷与答题纸一并收回。
    \item 回答各个问题时,将答案写在答题纸上,写在本试卷上无效。
\end{enumerate}

\vspace{0.5cm}
\noindent\textbf{一、选择题（30 pts.）}

\begin{enumerate}[label=\arabic*., itemsep=1.5ex]
    \item 电解时, 在阳极上首先发生氧化作用而放电的是: ( )
    \begin{tasks}(1)
        \task 标准还原电势最大者
        \task 标准还原电势最小者
        \task 考虑极化后, 实际上的不可逆还原电势最大者
        \task 考虑极化后, 实际上的不可逆还原电势最小者
    \end{tasks}
    \item 恒温恒压条件下, 某化学反应若在电池中可逆进行时吸热, 据此可以判断下列热力学量中何者一定大于零? ( )
    \begin{tasks}(4)
        \task $\Delta U$
        \task $\Delta H$
        \task $\Delta S$
        \task $\Delta G$
    \end{tasks}
    \item 对亚铁氧化铜负溶胶而言, 电解质 $\ce{KCl, CaCl_{2}, K_{2}SO_{4}, CaSO_{4}}$ 的聚沉能力顺序为: ( )
    \begin{tasks}(1)
        \task $KCl > CaCl_{2} > K_{2}SO_{4} > CaSO_{4}$
        \task $CaSO_{4} > CaCl_{2} > K_{2}SO_{4} > KCl$
        \task $CaCl_{2} > CaSO_{4} > KCl > K_{2}SO_{4}$
        \task $K_{2}SO_{4} > CaSO_{4} > CaCl_{2} > KCl$
    \end{tasks}
    \item
    \begin{enumerate}[label=(\arabic*), noitemsep, topsep=0pt, partopsep=0pt]
        \item 处于标准态的 $\ce{CO_{2}(g)}$ 和 $\ce{O_{2}(g)}$, 其标准燃烧焓值为零,
        \item 因为 $\Delta_\mathrm{r}G_\mathrm{m}^\ominus = -RT\ln K_p^\ominus$, 而 $K_p^\ominus$ 是由平衡时的组成表示的, 所以$\Delta_\mathrm{r}G_\mathrm{m}^\ominus$ 表示平衡时产物的吉布斯自由能与反应物的吉布斯自由能之差
        \item 水在 $25^\circ\mathrm{C}$ $p^\ominus$ 下蒸发, 求算熵变的公式为 $\Delta S_m^\ominus = (\Delta H_m^\ominus - \Delta G_m^\ominus)/T$
        \item 在恒温, 恒压下可逆电池反应, 求算熵变的公式为 $\Delta_\mathrm{r}S_\mathrm{m} = \Delta_\mathrm{r}H_\mathrm{m}/T$
    \end{enumerate}
    上述说法正确的是: ( )
    \begin{tasks}(4)
        \task 1, 2
        \task 2, 3
        \task 1, 3
        \task 3, 4
    \end{tasks}
    \item 电池在恒温、恒压及可逆情况下放电, 则其与环境的热交换为 ( )
    \begin{tasks}(2)
        \task $\Delta_\mathrm{r}H$
        \task $T\Delta_\mathrm{r}S$
        \task 一定为零
        \task 与 $\Delta_\mathrm{r}H$ 及 $T\Delta_\mathrm{r}S$ 均无关
    \end{tasks}
    \item 苯在一个刚性的绝热容器中燃烧, 则: ( )
    \[
    \ce{C_{6}H_{6}(l) + (15/2)O_{2}(g) \rightleftharpoons 6 CO_{2} + 3H_{2}O(g)}
    \]
    \begin{tasks}(2)
        \task $\Delta U = 0$, $\Delta H < 0$, $Q = 0$
        \task $\Delta U = 0$, $\Delta H > 0$, $W = 0$
        \task $\Delta U = 0$, $\Delta H = 0$, $Q = 0$
        \task $\Delta U \neq 0$, $\Delta H \neq 0$, $Q = 0$
    \end{tasks}
    \item 一可逆热机与另一不可逆热机在其他条件都相同时, 燃烧等量的燃料, 则可逆热机牵引的列车行走的距离: ( )
    \begin{tasks}(4)
        \task 较长
        \task 较短
        \task 一样
        \task 不一定
    \end{tasks}
    \item 对于一定量的理想气体, 下列过程可能发生的是: ( )
    \begin{enumerate}[label=(\arabic*), noitemsep, topsep=0pt, partopsep=0pt]
        \item 对外作功, 同时放热
        \item 体积不变, 而温度上升, 并且是绝热过程, 无非体积功
        \item 恒压下温度不变
        \item 恒温下绝热膨胀做功
    \end{enumerate}
    \begin{tasks}(4)
        \task (1), (4)
        \task (2), (3)
        \task (3), (4)
        \task (1), (2)
    \end{tasks}
    \item $298\ \mathrm{K}$、$0.1\ \mathrm{mol\cdot dm^{-3}}$ 的 $\ce{HCl}$ 溶液中, 氢电极的热力学电势为 $-0.06\ \mathrm{V}$, 电解此溶液时, 氢在铜电极上的析出电势 $\ce{H_{2}}$ 为: ( )
    \begin{tasks}(4)
        \task 大于 $-0.06\ \mathrm{V}$
        \task 等于 $-0.06\ \mathrm{V}$
        \task 小于 $-0.06\ \mathrm{V}$
        \task 不能判定
    \end{tasks}
    \item 斜方硫的燃烧热等于 ( )
    \begin{tasks}(2)
        \task $\ce{SO_{2}(g)}$ 的生成热
        \task $\ce{SO_{3}(g)}$ 的生成热
        \task 单斜硫的燃烧热
        \task 零
    \end{tasks}
    \item 在电池中, 当电池反应达到平衡时, 电池的电动势等于: ( )
    \begin{tasks}(4)
        \task 标准电动势
        \task $\dfrac{RT}{zF}\ln K^\ominus$
        \task 零
        \task 不确定
    \end{tasks}
    \item 已知 $\varphi^\ominus_{\ce{Fe^{3+}/Fe^{2+}}} = 0.77\ \mathrm{V}$, $\varphi^\ominus_{\ce{Sn^{4+}/Sn^{2+}}} = 0.15\ \mathrm{V}$ 
    
    以$\ce{2Fe^{3+} + Sn^{2+} = Sn^{4+} + 2Fe^{2+}}$ 反应组成原电池, 其标准电动势 $E^\ominus$ 值为: ( )
    \begin{tasks}(4)
        \task $1.09\ \mathrm{V}$
        \task $0.92\ \mathrm{V}$
        \task $1.39\ \mathrm{V}$
        \task $0.62\ \mathrm{V}$
    \end{tasks}
    \item $0.1\ \mathrm{mol\cdot kg^{-1}}$ 氯化钡水溶液的离子强度为: ( )
    \begin{tasks}(4)
        \task $0.1\ \mathrm{mol\cdot kg^{-1}}$
        \task $0.15\ \mathrm{mol\cdot kg^{-1}}$
        \task $0.2\ \mathrm{mol\cdot kg^{-1}}$
        \task $0.3\ \mathrm{mol\cdot kg^{-1}}$
    \end{tasks}
    \item 某反应的反应物消耗一半的时间正好是反应物消耗 $1/4$ 的时间的 $2$ 倍, 则该反应的级数是: ( )
    \begin{tasks}(4)
        \task $0.5$ 级反应
        \task $0$ 级反应
        \task $1$ 级反应
        \task $2$ 级反应
    \end{tasks}
    \item 已知$\varphi^\ominus_{\ce{Cl_{2}/Cl^{-}}} = 1.36\ \mathrm{V}$, $\eta_{\ce{(Cl_{2})}} = 0\ \mathrm{V}$, $\varphi^\ominus_{\ce{O_{2}/OH^{-}}} = 0.401\ \mathrm{V}$, $\eta_{\ce{(O_{2})}} = 0.8\ \mathrm{V}$. 以石墨为阳极, 电解 $0.01\ \mathrm{mol\cdot kg^{-1}}$ $\ce{NaCl}$ 溶液, 在阳极上首先析出: ( )
    \begin{tasks}(2)
        \task $\ce{Cl_{2}}$
        \task $\ce{O_{2}}$
        \task $\ce{Cl_{2}}$ 与 $\ce{O_{2}}$ 混合气体
        \task 无气体析出
    \end{tasks}
\end{enumerate}

\vspace{0.5cm}
\noindent\textbf{二、计算题（15 pts.）}

在 $630\ \mathrm{K}$ 时, 反应: $\ce{2HgO(s) <=> 2Hg(g) + O_{2}(g)}$ 的 $\Delta_\mathrm{r}G_\mathrm{m}^\ominus = 44.3\ \mathrm{kJ\cdot mol^{-1}}$.
\begin{enumerate}[label=(\arabic*), itemsep=1ex]
    \item 求上述反应的标准平衡常数 $K_p^\ominus$
    \item 求 $630\ \mathrm{K}$ 时 $\ce{HgO(s)}$ 的分解压力;
    \item 若将 $\ce{HgO(s)}$ 投入到 $630\ \mathrm{K}$, $1.013\times10^5\ \mathrm{Pa}$ 的纯 $\ce{O_{2}}$ 气的定体积容器中, 在 $630\ \mathrm{K}$ 时使其达平衡, 求与 $\ce{HgO(s)}$ 呈平衡的气相中 $\ce{Hg(g)}$ 的分压力.
\end{enumerate}

\vspace{0.5cm}
\noindent\textbf{三、计算题（15 pts.）}

$25^\circ\mathrm{C}$ 时甲烷溶在苯中, 当平衡浓度 $x_{\ce{(CH_{4})}} = 0.0043$ 时, $\ce{CH_{4}}$ 在平衡气相中的分压为 $245\ \mathrm{kPa}$. 已知 $25^\circ\mathrm{C}$ 时液态苯和苯蒸气的标准生成焓分别为 $48.66$ 和 $82.93\ \mathrm{kJ\cdot mol^{-1}}$, 苯在 $101\ 325\ \mathrm{Pa}$ 下的沸点为 $80.1^\circ\mathrm{C}$. 试计算
\begin{enumerate}[label=(\arabic*), itemsep=1ex]
    \item $25^\circ\mathrm{C}$ 时当 $x_{\ce{(CH_{4})}} = 0.01$ 时的甲烷苯溶液的蒸气总压 $p$ .
    \item 与上述溶液成平衡的气相组成 $y_{\ce{(CH_{4})}}$ .
\end{enumerate}

\vspace{0.5cm}
\noindent\textbf{四、计算题（15 pts.）}

$\ce{Mg}$ ( 熔点为 $648^\circ\mathrm{C}$ ) 和 $\ce{Cu}$ ( 熔点为 $1085^\circ\mathrm{C}$ ) 形成两种化合物: $\ce{MgCu_{2}}$ ( 熔点为 $800\ ^\circ\mathrm{C}$ ), $\ce{Mg_{2}Cu}$ ( 熔点为 $580^\circ\mathrm{C}$ ); 形成三个低共熔混合物, 其组成分别为 $w_\mathrm{Mg} = 0.10, w_\mathrm{Mg} = 0.33, w_\mathrm{Mg} = 0.65$, 它们的熔点分别为 $690^\circ\mathrm{C}, 560^\circ\mathrm{C}$ 和 $380^\circ\mathrm{C}$ 
\begin{enumerate}[label=(\arabic*), itemsep=1ex]
    \item 利用上述数据绘制 $\ce{Cu-Mg}$ 体系的相图.
    \item 画出组成为 $w_\mathrm{Mg} = 0.25$ 的 $900^\circ\mathrm{C}$ 的熔体降温到 $100^\circ\mathrm{C}$ 时的步冷曲线.
    \item 当 (2) 中的熔体 $1\ \mathrm{kg}$ 降温至无限接近 $560^\circ\mathrm{C}$ 时, 得到何种化合物? 其质量几何?
\end{enumerate}

\vspace{0.5cm}
\noindent\textbf{五、计算题（20 pts.）}

一直到 $1000\ p^\ominus$, 氮气仍服从下列状态方程式: $pV_\mathrm{m} = RT + bp$, 式中常数 $b = 3.90\times10^{-2}\ \mathrm{dm^3\cdot mol^{-1}}$ . 在 $500\ \mathrm{K}$, $1\ \mathrm{mol}\ \ce{N_{2}(g)}$ 从 $1000\ p^\ominus$ 等温膨胀到 $p^\ominus$ 计算 $U_\mathrm{m}, \Delta V_\mathrm{m}, \Delta G_\mathrm{m}, \Delta A_\mathrm{m}$ 及 $S_\mathrm{m}$ .

\vspace{0.5cm}
\noindent\textbf{六、计算题（20 pts.）}

反应 $\ce{[Co(NH_{3})_{5}F]^{2+} + H_{2}O \to [Co(NH_{3})_{5}H_{2}O]^{3+} + F^{-}}$ 是一个酸催化反应, 设反应的速率公式可表示为:
\[
-\frac{\mathrm{d}[\ce{Co(NH_{3})_{5}F^{2+}}]}{\mathrm{d}t} = k[\ce{Co(NH_{3})_{5}F^{2+}}]^a[\ce{H^{+}}]^b
\]
今摘取关于此反应的三次实验结果列于下表:

\begin{center}
\begin{tabular}{ccccc}
$[\ce{Co(NH_{3})_{5}F^{2+}}]/\mathrm{mol\cdot dm^{-3}}$ & $[\ce{H^{+}}]/\mathrm{mol\cdot dm^{-3}}$ & $t/^\circ\mathrm{C}$ & $t_{1/2}/\mathrm{h}$ & $t_{3/4}/\mathrm{h}$ \\
\hline
$0.1$ & $0.01$ & $25$ & $1$ & $2$ \\
$0.2$ & $0.02$ & $25$ & $0.5$ & $1$ \\
$0.1$ & $0.01$ & $25$ & $0.5$ & $1$ \\
\end{tabular}
\end{center}

\begin{enumerate}[label=(\arabic*), itemsep=1ex]
    \item 试求 $a, b$ .
    \item 试计算速变常数 $k$ 值
    \item 求该反应的活化能 $E_\mathrm{a}$
\end{enumerate}

\vspace{0.5cm}
\noindent\textbf{七、计算题（20 pts.）}

电池 $\ce{Cd | Cd(OH)_{2} | NaOH(0.01\ mol\cdot kg^{-1}) | H_{2}(p^\ominus) | Pt}$ $298\ \mathrm{K}$ 时电动势为 $E = 0.000\ \mathrm{V}$, $(\partial E/\partial T)_p = 0.002\ \mathrm{V\cdot K^{-1}}$, $\varphi^\ominus_{\ce{(Cd^{2+}/Cd)}} = -0.403\ \mathrm{V}$ .
\begin{enumerate}[label=(\arabic*), itemsep=1ex]
    \item 写出两电极反应和电池反应
    \item 求电池反应的 $\Delta_\mathrm{r}G_\mathrm{m}$、$\Delta_\mathrm{r}H_\mathrm{m}$、$\Delta_\mathrm{r}S_\mathrm{m}$
    \item 求 $\ce{Cd(OH)_{2}}$ 的溶度积常数 $K_\mathrm{sp}$
\end{enumerate}

\vspace{0.5cm}
\noindent\textbf{八、计算题（15 pts.）}

电极 $\ce{Ag^{+} | Ag}$ 和 $\ce{Cr^{-} | AgCl(s) | Ag}$ 在 $298\ \mathrm{K}$ 时的标准电极电位分别为 $0.7991\ \mathrm{V}$ 和 $0.2224\ \mathrm{V}$ .
\begin{enumerate}[label=(\arabic*), itemsep=1ex]
    \item 计算 $\ce{AgCl}$ 在水中的饱和溶液的浓度
    \item 用德拜-休克尔方程式计算 $298\ \mathrm{K}$ 时 $\ce{AgCl}$ 在 $0.01\ \mathrm{mol\cdot kg^{-1}}\ \ce{KNO_{3}}$ 溶液中的溶解度; 已知 $A = 0.509(\mathrm{mol\cdot kg^{-1}})$
    \item 计算反应 $\ce{AgCl(s) <=> Ag^{+}(aq) + Cl^{-}(aq)}$ 的 $\Delta_\mathrm{r}G_\mathrm{m}^\ominus$
\end{enumerate}